\section{Mean-field derivations}
\label{supp:s2-mean-field}

This note provides the derivations omitted from the main text for the mean-field critical point and the Ginzburg--Landau form.

\subsection{Mean-field update and collective response}

Under well-mixed assumptions, each update samples $k$ independent risk signals from the environment, each labeled ``risk'' with probability $p_{\text{env}}$. Let $S\sim\text{Binomial}(k,p_{\text{env}})$ and $\hat{p}=S/k$ be the perceived risk fraction. Given thresholds $\theta<\phi$, an updated individual moves to $H$ if $\hat{p}\ge \phi$, to $L$ if $\hat{p}\le \theta$, and otherwise remains in $M$. Therefore,
\begin{align}
P_H(p_{\text{env}}) &= \Pr(\hat{p}\ge \phi)=\sum_{s=\lceil k\phi\rceil}^{k}\binom{k}{s}p_{\text{env}}^{\,s}(1-p_{\text{env}})^{k-s},\\
P_L(p_{\text{env}}) &= \Pr(\hat{p}\le \theta)=\sum_{s=0}^{\lfloor k\theta\rfloor}\binom{k}{s}p_{\text{env}}^{\,s}(1-p_{\text{env}})^{k-s}.
\end{align}
The expected polarization direction after an update is the difference between high- and low-arousal outcomes,
\begin{equation}
\mathcal{S}(p_{\text{env}})=P_H(p_{\text{env}})-P_L(p_{\text{env}}).
\end{equation}
With asynchronous updates and a continuous-time approximation, the mean-field dynamics take the form
\begin{equation}
\frac{dq}{dt}=-q+\mathcal{S}(p_{\text{env}}),
\end{equation}
which is the mean-field form used in the main text.

\subsection{Linear stability and analytic critical point}

In the symmetric validation regime,
\begin{equation}
p^{\text{main}}(q)=\frac{1-q}{2},\qquad p^{\text{we}}(q)=\frac{1+q}{2}.
\end{equation}
The environmental risk is the weighted average of the two channels,
\begin{equation}
p_{\text{env}}(q;r)=\frac{(1-r)n_m p^{\text{main}}(q)+r n_w p^{\text{we}}(q)}{(1-r)n_m+r n_w},
\end{equation}
where $n_m$ and $n_w$ are baseline channel strengths and $r\in[0,1]$ controls their relative weight. Expanding around $q=0$ yields
\begin{equation}
p_{\text{env}}(q;r)=\frac{1}{2}+\Gamma(r)\,q+O(q^3),
\end{equation}
where the feedback gradient is
\begin{equation}
\Gamma(r)=\left.\frac{\partial p_{\text{env}}}{\partial q}\right|_{q=0}
=\frac{r n_w-(1-r)n_m}{2\left((1-r)n_m+r n_w\right)}.
\end{equation}
Next, expand the collective response around neutrality $p=1/2$:
\begin{equation}
\mathcal{S}(p)=\mathcal{S}(1/2)+\chi\left(p-\frac{1}{2}\right)+O\!\left(\left(p-\frac{1}{2}\right)^3\right),\qquad
\chi=\left.\frac{d\mathcal{S}}{dp}\right|_{p=1/2}.
\end{equation}
Under symmetry, $\mathcal{S}(1/2)=0$, so substituting $p_{\text{env}}(q;r)$ into the mean-field dynamics gives, to leading order,
\begin{equation}
\frac{dq}{dt}=(-1+\chi\Gamma)\,q+O(q^3).
\end{equation}
Therefore the neutral fixed point $q=0$ loses stability when the linear restoring coefficient changes sign:
\begin{equation}
\chi\Gamma(r_c)=1.
\end{equation}
Solving for $r$ yields the closed-form critical point reported in the main text:
\begin{equation}
r_c=\frac{n_m(\chi+2)}{n_m(\chi+2)+n_w(\chi-2)}.
\end{equation}
When $\chi\le 2$, the denominator does not admit a critical point in $[0,1]$, consistent with the absence of a transition for weak psychological sensitivity.

\subsection{Ginzburg--Landau form and noise term}

Keeping the next nonlinearity in the expansion yields a cubic term. Writing
\begin{equation}
\mathcal{S}(p)=\chi\left(p-\frac{1}{2}\right)+\frac{\mathcal{S}^{(3)}(1/2)}{6}\left(p-\frac{1}{2}\right)^3+O\!\left(\left(p-\frac{1}{2}\right)^5\right),
\end{equation}
and using $p_{\text{env}}-1/2\approx \Gamma q$, we obtain
\begin{equation}
\frac{dq}{dt}=(-1+\chi\Gamma)\,q+\frac{\mathcal{S}^{(3)}(1/2)}{6}\Gamma^3 q^3+\eta(t)+\cdots
=-\alpha q-u q^3+\eta(t)+\cdots,
\end{equation}
where $\alpha=1-\chi\Gamma$ and $u=-\mathcal{S}^{(3)}(1/2)\Gamma^3/6$. For the parameter ranges considered, $u>0$, so the dynamics can be written in the effective-potential form
\begin{equation}
\frac{dq}{dt}=-\frac{\partial V_{\text{eff}}(q)}{\partial q}+\eta(t),\qquad
V_{\text{eff}}(q)=\frac{1}{2}\alpha q^2+\frac{u}{4}q^4,
\end{equation}
which gives a pitchfork bifurcation when $\alpha$ changes sign.

The noise term $\eta(t)$ represents stochastic fluctuations induced by finite population size and discrete updates. In the large-$N$ limit, these fluctuations can be approximated as Gaussian white noise with an amplitude that scales as $1/\sqrt{N}$, implying a variance scaling as $1/N$.
